\subsection{The real condition in the pro-two tower is necessary}

\begin{result}{E.157}{Known}{The actual real pro-two section classes of the tripod}
The geometrically pro-two arithmetic fundamental group of
\(U_{\mathbb R}=\mathbb P^1_{\mathbb R}\setminus\{0,1,\infty\}\)
is
\[
 E^{(2)}_{U_{\mathbb R}}\simeq
 C_2\amalg_{\mathrm{pro}\text{-}2}
 C_2\amalg_{\mathrm{pro}\text{-}2}C_2.
 \tag{E.157.1}
\]
Its geometric kernel is free pro-two of rank two, and its
three geometric-conjugacy classes of sections are precisely
the ordinary real point classes of the intervals
\((-\infty,0)\), \((0,1)\), and \((1,\infty)\).
They are distinguished by the signs of \(t,1-t\).
The two real tangential classes at each cusp give the two
adjacent interval classes.

Consequently any full global geometrically pro-two section
has an ordinary real point class. For a fixed ordinary
completed pair, requiring this particular real class at
every nilpotent level loses no full pro-two survivor.
Together with A.21, A.31 and E.155, its survival is equivalent
to nonemptiness at every level of the compact section fibres
with the prescribed finite boundary matches and that real
class. At each such matched predecessor the global lifting
obstruction vanishes; the remaining transition test is
simultaneous first-cohomology matching.
\end{result}
\begin{proof}
This is a direct application of a published pro-two fixed-point
theorem, with the comparison and hypotheses made explicit.
Take the real base point \(b=1/2\). Choose the two based
topological meridians \(x,y\) around \(0,1\), respectively,
with real stems. Complex conjugation \(c\) fixes the stems
and reverses the small circles. Thus its actual based action
is
\[
                         c(x)=x^{-1},\qquad c(y)=y^{-1}.
 \tag{E.157.2}
\]
The Riemann existence equivalence
\citep[XII, Theorem~5.1, recomposed p.~251]{RootSGA1}
identifies finite algebraic and finite topological covers
and is compatible with this based conjugation. Since \(b\)
supplies the real splitting, the actual arithmetic extension
is the corresponding semidirect product. Passing to the
geometrically pro-two quotient gives
\[
 F^{(2)}(x,y)\rtimes\langle c\rangle
 =\langle x,y,c\mid c^2=1,\ (xc)^2=1,\ (yc)^2=1
                                      \rangle_{\mathrm{pro}\text{-}2}.
 \tag{E.157.3}
\]
It is itself pro-two, and has the universal property of the
maximal pro-two quotient of the full arithmetic group:
any finite two-group quotient restricts on the geometric
group through its pro-two completion. Setting
\(s_0=c,s_1=xc,s_2=yc\) gives mutually inverse presentations
of (E.157.3) and the free pro-two product of three groups
of order two. This proves (E.157.1), including the actual
arithmetic action rather than only an outer action.

To classify its sections, let \(X\) be the finite graph
formed from two circles with one common point. Reflect
each circle, fixing the common point. Subdivide at its
three fixed points; this is a finite \(C_2\)-CW complex.
Its based fundamental group and involution are exactly
the free group on \(x,y\) with (E.157.2), so its pro-two
arithmetic extension is (E.157.3).

The graph is \((\{2\},2)\)-erasable in the terminology of
\citet[Definition~1.2, author p.~2]{CohPal2011}.
Indeed its cohomology with any local coefficients vanishes
in degrees at least two. A degree-one class in a finite
two-primary local system with two-group monodromy is a
cocycle into its finite module \(M\). The corresponding
homomorphism to \(M\rtimes\operatorname{im}(\rho)\) has
finite two-group image; its kernel defines a finite
normal two-cover on which the monodromy and that class
are both zero. This verifies the required erasability
also if \emph{two-torsion} is read as killed by two.

Apply \citet[Theorem~2.4, author p.~5; Definition~2.3,
author pp.~4--5]{CohPal2011} with \(p=2\) and
\(\Sigma=\{2\}\). It identifies sections modulo conjugation
by the geometric pro-two group with the connected
components of \(X^{C_2}\). There are exactly three.
The fixed common point gives \(c\), and the opposite
fixed point of each circle gives \(xc\) or \(yc\), up
to geometric conjugation. For example, a half-circle
path to the opposite point followed by the inverse of
its conjugate is the full meridian.
The same based paths in \(U(\mathbb C)\) identify these
classes with points in the three real intervals. Their
abelianized classes are respectively
\[
                    (0,0),\qquad(1,0),\qquad(0,1)
                    \quad\hbox{in }\mathbb F_2^2,
 \tag{E.157.4}
\]
the three ordinary sign patterns. This also verifies
their distinctness in the actual Kummer coordinates.
We have used the topological pro-\(\{2\}\) theorem,
not inferred its assertion from a full-profinite real
section theorem.

For the cusp assertion, a real tangential parameter has
two completed Kummer classes, its two signs. At cusp
zero the pair has signs \((\pm,+)\); at cusp one it
has signs \((+,\pm)\); and at infinity it has signs
\((\pm,\mp)\). Each is one of (E.157.4), so the preceding
classification gives exactly the adjacent intervals.
No additional real class remains merely because its
construction used a cusp.

Restriction of a full global pro-two section is a section
of (E.157.3), proving necessity. The signs fix its interval
class once the global abelian pair is fixed. In the case
of an ordinary completed pair, the forbidden sign
pattern \((-,-)\) is excluded already by the real
component of the cup calculation A.25.
At every finite support prime \(q\ne2\), A.31 makes the
full local class the specified boundary point class.
Hence a full survivor belongs at every finite depth to
the stated matched fibres.

Conversely, those fibres are compact by A.20--A.21:
one fixes the global abelian cocycle, and each local
matching condition is closed because it is the compact
geometric conjugacy orbit of a specified point class.
Their transition maps preserve all the conditions.
Nonemptiness at every level therefore gives a compatible
inverse-limit section by the finite-intersection argument
of A.21; surjectivity of transition maps is not assumed.
E.155 supplies the final lifting assertion. It does not
make the subsequent simultaneous matching equations
automatically consistent.
\end{proof}
